\documentclass[../BTL.tex]{subfiles}
\begin{document}

\section{Kafka}
Apache Kafka là hệ thống xử lý dữ liệu phân tán hoạt động theo mô hình luồng sự kiện (event streaming), cho phép truyền tải và xử lý dữ liệu thời gian thực với độ trễ thấp và thông lượng cao. Kafka sử dụng thuật toán Raft \cite{kafka_internals} để quản lý metadata và bầu chọn leader cho mỗi phân vùng (partition) thông qua cơ chế quorum. Kafka dùng cơ chế phân vùng (partition) và sao chép (replication) để phân tán dữ liệu trên nhiều broker, tăng khả năng chịu lỗi và đạt thông lượng cao \cite{kafka_design}.

\subsection{Kiến trúc cơ bản của Kafka}
Kiến trúc của Kafka được tổ chức theo mô hình phân tán với các thành phần chính sau \cite{kafka_official,streaming_systems}:

\begin{itemize}
    \item \textbf{Broker:} Máy chủ trong cụm chịu trách nhiệm lưu trữ và phục vụ dữ liệu. Mỗi broker quản lý một tập các partition thuộc các topic khác nhau.

    \item \textbf{Topic:} Đơn vị phân loại dữ liệu trong Kafka. Mỗi topic đại diện cho một luồng dữ liệu có thứ tự. Producer gửi bản ghi tới topic và consumer đọc bản ghi từ topic.

    \item \textbf{Partition:} Mỗi topic được chia thành nhiều partition để cho phép song song hóa việc đọc và ghi dữ liệu. Các bản ghi trong mỗi partition được gán một số offset (độ dời) duy nhất. Consumer sử dụng offset để theo dõi vị trí đọc.

    \item \textbf{Producer / Consumer:} Producer gửi bản ghi tới các topic và phân phối tới các partition dựa trên khóa phân vùng. Consumer đăng ký theo dõi topic và liên tục kéo (pull) các bản ghi mới từ broker. Kafka đảm bảo mỗi bản ghi trong một partition chỉ được đọc bởi một consumer trong cùng một consumer group.

    \item \textbf{Replication:} Mỗi partition có một bản sao chính (leader) và nhiều bản sao dự phòng (follower). Tất cả thao tác ghi và đọc đều qua leader, và leader đồng bộ hóa dữ liệu tới các follower. Kafka hỗ trợ cấu hình mức độ sao chép (replication factor) và ba mức xác nhận: \textit{acks=0}, \textit{acks=1}, và \textit{acks=all}.

    \item \textbf{KRaft:} Cơ chế quản lý metadata nội bộ của Kafka, thay thế cho ZooKeeper từ phiên bản 3.4. Một broker đóng vai trò controller, chịu trách nhiệm quản lý trạng thái các broker và xử lý sự kiện thay đổi trạng thái \cite{kafka_internals}.
\end{itemize}

Kafka hoạt động theo mô hình subscribe-publish (xuất bản - đăng ký), giúp tách rời producer và consumer, cho phép hệ thống mở rộng linh hoạt. Dữ liệu được lưu trữ trong các partition theo cấu trúc log có thể cấu hình (retention period), cho phép consumer đọc lại dữ liệu cũ thông qua điều chỉnh offset \cite{kafka_definitive, streaming_systems}.

Ưu điểm của Kafka: khả năng chịu tải và mở rộng cao, đảm bảo tính nhất quán (log-based consistency), độ tin cậy cao (nhờ replication), xử lý dòng sự kiện thời gian thực, và lưu trữ bền vững cho phép tái xử lý.

Nhược điểm của Kafka: phức tạp trong triển khai và quản lý, yêu cầu tài nguyên đáng kể, khả năng quản lý dữ liệu cũ hạn chế, và rủi ro quá tải từ sự kiện không cần thiết \cite{kafka_official}.


\section{Apache Spark}
Apache Spark là framework xử lý dữ liệu phân tán mã nguồn mở, sử dụng mô hình xử lý trong bộ nhớ (in-memory computing) để đạt hiệu suất cao hơn so với các framework dựa trên đĩa truyền thống \cite{spark_rdd, spark_definitive}. Spark được xây dựng trên khái niệm Resilient Distributed Datasets (RDD) - tập dữ liệu phân tán có khả năng phục hồi lỗi (fault tolerance), từ đó phát triển thêm các API trừu tượng ở mức cao hơn như DataFrame và Dataset \cite{spark_rdd, spark_sql}.

\subsection{Các thành phần chính của Spark}
Spark được tổ chức thành các tầng thành phần xếp chồng lên nhau \cite{spark_definitive, spark_official}:

\begin{itemize}
    \item \textbf{Spark Core:} Nền tảng của toàn bộ hệ thống, chịu trách nhiệm quản lý bộ nhớ, lập lịch tác vụ (task scheduling), theo dõi lỗi (fault recovery), và tương tác với hệ thống lưu trữ bên ngoài. Cung cấp API cho Java, Scala, Python và R \cite{spark_rdd}.

    \item \textbf{Spark SQL:} Cung cấp kiểu trừu tượng DataFrame/Dataset, hỗ trợ dữ liệu có cấu trúc và nửa cấu trúc. Sử dụng bộ tối ưu hóa Catalyst để phân tích, tối ưu và thực thi truy vấn tự động \cite{spark_sql}.

    \item \textbf{Spark Streaming / Structured Streaming:} Xử lý luồng dữ liệu bằng cách chia luồng thành các micro-batch và áp dụng RDD transformation lên từng lô nhỏ. Structured Streaming mở rộng DataFrame để xử lý luồng theo cách khai báo (declarative), tự động quản lý checkpoint và offset \cite{spark_streaming, spark_structured_streaming}.

    \item \textbf{MLlib:} Thư viện học máy tích hợp sẵn của Spark, cung cấp các thuật toán phân loại (classification), hồi quy (regression), gom cụm (clustering), lọc cộng tác (collaborative filtering) và giảm chiều (dimensionality reduction) \cite{spark_mllib}.

    \item \textbf{GraphX:} Khung xử lý đồ thị phân tán trên nền tảng Spark, cung cấp bộ API linh hoạt và các thuật toán đồ thị phân tán phổ biến như PageRank, connected components và triangle counting \cite{spark_graphx}.
\end{itemize}

Ưu điểm của Spark: tương thích cao với các nguồn dữ liệu Hadoop và nhiều hệ thống khác (Kafka, Cassandra, JDBC), hỗ trợ nhiều ngôn ngữ (Scala, Java, Python, R), xử lý trong bộ nhớ cho tốc độ cao, hỗ trợ xử lý theo lô (batch processing), phân tích thời gian thực, học máy và xử lý đồ thị trong một framework duy nhất, và hệ sinh thái phong phú \cite{spark_definitive, spark_official}.

\subsection{Cơ chế hoạt động của Spark}
Spark giải quyết các hạn chế của Hadoop MapReduce bằng cách tối ưu hóa quy trình xử lý dữ liệu \cite{spark_rdd, spark_definitive}:

\begin{itemize}
    \item \textbf{Hadoop MapReduce} yêu cầu đọc-ghi đĩa nhiều lần qua HDFS cho mỗi giai đoạn (stage), gây độ trễ I/O đáng kể, đặc biệt khi cần nhiều vòng lặp tính toán liên tiếp \cite{hadoop_mapreduce}.

    \item \textbf{Xử lý trong bộ nhớ (in-memory processing):} Spark đọc dữ liệu một lần vào bộ nhớ, thực hiện các thao tác biến đổi (transformation) trực tiếp trên dữ liệu trong bộ nhớ, và chỉ ghi kết quả cuối cùng ra đĩa khi cần \cite{spark_rdd}.

    \item \textbf{Tái sử dụng dữ liệu (data reuse):} Spark sử dụng bộ nhớ đệm (cache) để lưu dữ liệu đã xử lý, cho phép nhiều thao tác tái sử dụng mà không cần tính toán lại. Đặc biệt hữu ích cho thuật toán học máy lặp và xử lý đồ thị \cite{spark_rdd, spark_sql}.

    \item \textbf{Đồ thị xử lý có hướng không chu trình (DAG):} Spark Core chuyển chương trình thành DAG gồm các giai đoạn (stage) và tác vụ (task). Mỗi giai đoạn chứa các tác vụ thực thi song song trên các executor, và Spark tối ưu hóa DAG bằng cách gộp các thao tác có thể thực hiện cùng nhau \cite{spark_rdd}.

    \item \textbf{Đánh giá lười (lazy evaluation):} Các phép biến đổi (transformation) trên RDD hoặc DataFrame không được thực thi ngay, chỉ được ghi nhận trong DAG. Khi gặp thao tác hành động (action) như collect, count hoặc save, Spark mới thực sự tính toán toàn bộ DAG. Điều này cho phép Spark tối ưu kế hoạch thực thi toàn cục \cite{spark_sql}.
\end{itemize}

Spark hỗ trợ bốn chế độ triển khai: Spark Standalone, Apache YARN, Mesos, và Kubernetes \cite{spark_official}.


\section{Redis}
Redis là hệ thống lưu trữ dữ liệu dạng key-value trong bộ nhớ. Nhờ lưu trữ toàn bộ dữ liệu trên RAM và thiết kế đơn luồng (single-threaded), Redis đạt được độ trễ rất thấp cho cả thao tác đọc và ghi, phù hợp với các trường hợp sử dụng đòi hỏi phản hồi nhanh \cite{redis_official,redis_wikipedia}.

\subsection{Các kiểu dữ liệu và tính năng chính}
Redis hỗ trợ nhiều kiểu dữ liệu phức tạp ngoài string thuần túy, cho phép thao tác trực tiếp trên cấu trúc dữ liệu phía server thay vì chỉ trả về giá trị đơn thuần \cite{redis_official,redis_definitive}:

\begin{itemize}
    \item \textbf{String:} Kiểu cơ bản nhất, lưu trữ chuỗi byte có thể là văn bản, JSON, hoặc số.

    \item \textbf{Hash:} Bảng ánh xạ giữa các trường (field) và giá trị, phù hợp để biểu diễn đối tượng \cite{redis_official}.

    \item \textbf{List:} Danh sách các chuỗi có thứ tự, hỗ trợ thao tác hai đầu (push/pop) với độ phức tạp O(1).

    \item \textbf{Set:} Tập hợp các phần tử không trùng lặp, hỗ trợ các phép toán hợp (union), giao (intersection) và hiệu (difference) \cite{redis_definitive}.

    \item \textbf{Sorted Set (ZSet):} Tập hợp có thứ tự, mỗi phần tử gắn kèm một điểm số (score) dùng để sắp xếp, thường dùng cho bảng xếp hạng (leaderboard) \cite{redis_official}.

    \item \textbf{Stream:} Cấu trúc log cho phép lưu trữ nhiều trường và giá trị theo thứ tự thời gian tại một khóa duy nhất, hỗ trợ mô hình pub/sub và consumer group \cite{redis_wikipedia,redis_official}.

    \item \textbf{JSON:} Redis hỗ trợ kiểu dữ liệu JSON cho phép lưu trữ và truy vấn tài liệu JSON trực tiếp \cite{redis_official}.

    \item \textbf{Vector:} Redis hỗ trợ vector embedding phục vụ ứng dụng AI và tìm kiếm ngữ nghĩa (semantic search) \cite{redis_official}.
\end{itemize}

Ngoài các kiểu dữ liệu, Redis còn cung cấp các tính năng quan trọng: Pub/Sub cho mô hình nhắn tin thời gian thực, Transaction cho phép thực thi nhóm lệnh nguyên tử, và Lua scripting cho phép chạy script tùy chỉnh trên server \cite{redis_official,redis_wikipedia}.

\subsection{Persistence và Replication}
Mặc dù dữ liệu được lưu trữ trong RAM, Redis hỗ trợ hai cơ chế ghi dữ liệu ra đĩa để đảm bảo tính bền vững \cite{redis_wikipedia}: snapshotting (RDB) - chụp trạng thái toàn bộ dữ liệu tại một thời điểm theo định kỳ; và journaling (AOF - Append Only File) - ghi lại mỗi thao tác sửa đổi vào file log, cho phép khôi phục chính xác hơn. Mặc định, Redis ghi dữ liệu xuống đĩa ít nhất mỗi hai giây, trong trường hợp sự cố hệ thống, chỉ mất vài giây dữ liệu \cite{redis_wikipedia}.

Redis hỗ trợ master-replica replication, cho phép dữ liệu từ một Redis server sao chép tới nhiều replica. Tính năng này phục vụ cho việc mở rộng khả năng đọc và đảm bảo dự phòng dữ liệu \cite{redis_wikipedia}. Redis Cluster cho phép phân tán dữ liệu trên nhiều node, có thể mở rộng tới 1.000 node và duy trì hoạt động khi một số node bị lỗi \cite{redis_wikipedia,redis_official}.

Ưu điểm của Redis: độ trễ cực thấp (sub-millisecond), hỗ trợ nhiều kiểu dữ liệu phức tạp, nhân bản (replication) và phân cụm (clustering) cho khả năng mở rộng, Pub/Sub cho nhắn tin thời gian thực.

Nhược điểm của Redis: do lưu trữ hoàn toàn trên RAM nên chi phí vận hành cao, không phù hợp với dữ liệu lớn vượt quá dung lượng bộ nhớ khả dụng, và yêu cầu chiến lược persistence phù hợp để tránh mất dữ liệu khi sự cố \cite{redis_wikipedia}.


\section{FastAPI}
FastAPI là web framework dành cho Python, được thiết kế để xây dựng các giao diện lập trình ứng dụng HTTP với hiệu suất cao và tốc độ phát triển nhanh \cite{fastapi_official}. Framework sử dụng Python type hints để khai báo tham số, request body và response; tự động sinh tài liệu OpenAPI và cung cấp giao diện API tương tác (Swagger UI tại /docs và ReDoc tại /redoc) mà không cần cấu hình thêm \cite{fastapi_official,fastapi_wikipedia}. FastAPI được xây dựng trên Starlette cho phần web (ASGI) và Pydantic cho xác thực dữ liệu \cite{fastapi_official}.

FastAPI hỗ trợ đầy đủ các phương thức HTTP: GET, POST, PUT, DELETE, PATCH, OPTIONS. Ngoài ra, framework hỗ trợ WebSocket cho giao tiếp song công (full-duplex), Background Tasks cho xử lý bất đồng bộ sau khi trả response, và GraphQL thông qua tích hợp với Strawberry \cite{fastapi_official,fastapi_wikipedia}.

FastAPI dựa trên ASGI (Asynchronous Server Gateway Interface) - một chuẩn giao diện giữa ứng dụng Python và web server, cho phép xử lý bất đồng bộ \cite{fastapi_wikipedia}. Uvicorn đóng vai trò server mặc định, sử dụng mô hình đa tiến trình (multi-worker) với một tiến trình chính (main process) quản lý pool worker và phân phối request tới các worker \cite{fastapi_wikipedia}.

Kiến trúc của FastAPI tận dụng Pydantic để định nghĩa cấu trúc dữ liệu qua Python type hints. Khi một request đến, Pydantic tự động xác thực dữ liệu đầu vào, chuyển đổi từ JSON sang các kiểu Python tương ứng, và serialize dữ liệu đầu ra trở lại JSON - toàn bộ chỉ với một khai báo duy nhất \cite{fastapi_official}. Hệ thống tiêm phụ thuộc (Dependency Injection) cho phép khai báo các thành phần phụ thuộc (database session, authentication, logging) như tham số của endpoint; Framework tự động phân giải và tiêm các dependency này cho mỗi request, hỗ trợ cleanup sau khi request hoàn tất qua cơ chế yield \cite{fastapi_official,fastapi_wikipedia}. Hiệu suất của FastAPI đến từ kiến trúc bất đồng bộ (async) thuần túy và khả năng xử lý đồng thời nhiều request trên một vòng lặp sự kiện đơn luồng \cite{fastapi_wikipedia}.

Ưu điểm: hiệu suất cao, tự động sinh tài liệu OpenAPI, xác thực dữ liệu tự động qua Pydantic, hỗ trợ lập trình bất đồng bộ (async/await), tiêm phụ thuộc (Dependency Injection) mạnh mẽ, hỗ trợ WebSocket và background tasks.

Nhược điểm: hệ sinh thái bổ sung (third-party) còn đang phát triển, không phù hợp với các ứng dụng cần đồng bộ hóa phía server phức tạp như Django, và yêu cầu Python 3.8+ \cite{fastapi_official}.


\section{PostgreSQL}
PostgreSQL là hệ quản trị cơ sở dữ liệu quan hệ đối tượng (RDBMS) mã nguồn mở, nhấn mạnh tính mở rộng (extensibility) và tuân thủ chuẩn SQL cao \cite{postgresql_wikipedia}.

PostgreSQL quản lý concurrency (xử lý đồng thời) thông qua MVCC (Multiversion Concurrency Control), cho phép mỗi transaction có một "snapshot" riêng của database, loại bỏ phần lớn nhu cầu đọc khóa (read locks) \cite{postgresql_wikipedia}. Hệ thống cung cấp bốn mức độ cô lập transaction: Read Uncommitted, Read Committed, Repeatable Read và Serializable, và hỗ trợ serializability đầy đủ qua phương pháp SSI (Serializable Snapshot Isolation) \cite{postgresql_wikipedia}.

PostgreSQL hỗ trợ nhiều loại index phức tạp: B-tree, hash, GiST, GIN, SP-GiST và BRIN, cùng với khả năng tạo index trên biểu thức (expression index) và index một phần (partial index) \cite{postgresql_wikipedia,postgresql_official}. Về replication, PostgreSQL tích hợp sẵn binary replication dựa trên WAL (Write-Ahead Logging), cho phép cấu hình replication đồng bộ hoặc bất đồng bộ và phân tách truy vấn đọc (read scaling) trên các replica \cite{postgresql_wikipedia}.

PostgreSQL hỗ trợ nhiều kiểu dữ liệu phong phú bao gồm JSON/JSONB, array, geometric types, IPv4/IPv6, UUID, XML, và cho phép người dùng định nghĩa kiểu dữ liệu tùy chỉnh (user-defined types) cùng với các index method tương ứng \cite{postgresql_wikipedia,postgresql_official}. Tính năng Foreign Data Wrapper (FDW) cho phép PostgreSQL kết nối tới các nguồn dữ liệu bên ngoài như file hệ thống, RDBMS khác, hoặc dịch vụ web và truy vấn chúng như các bảng thông thường \cite{postgresql_wikipedia}.

Ưu điểm: mã nguồn mở với giấy phép PostgreSQL cho phép sử dụng tự do, tuân thủ chuẩn SQL cao, hỗ trợ nhiều kiểu dữ liệu phức tạp, MVCC giúp xử lý đồng thời hiệu quả mà không cần read locks, và khả năng mở rộng cao thông qua user-defined types và index methods \cite{postgresql_wikipedia,postgresql_official}.

Nhược điểm: hiệu suất trên các workload read-heavy có thể cần tuning đặc biệt khi dưới tải write nặng do đặc thù MVCC, và multi-master replication đồng bộ không được tích hợp sẵn trong core \cite{postgresql_wikipedia}.


\section{Streamlit}
Streamlit là framework mã nguồn mở viết bằng Python, được thiết kế để giúp các nhà khoa học dữ liệu và kỹ sư AI/ML nhanh chóng tạo ra các ứng dụng web dữ liệu chỉ với vài dòng code thuần Python, không yêu cầu kinh nghiệm về giao diện người dùng (front-end) \cite{streamlit_official}.

Nguyên tắc cốt lõi của Streamlit là khai báo widget (thành phần giao diện) tương đương với khai báo biến - không cần viết phần xử lý phía server (backend), định nghĩa tuyến đường (route), xử lý yêu cầu HTTP, hay viết HTML/CSS/JavaScript \cite{streamlit_official}. Mỗi khi file Python được lưu, ứng dụng tự động cập nhật ngay (tải lại tức thì - hot reload), cho phép lập trình viên thấy kết quả tức thì trong quá trình phát triển. API của Streamlit rất đơn giản: chỉ cần \texttt{import streamlit as st} và sử dụng các hàm như \texttt{st.write()}, \texttt{st.line\_chart()}, \texttt{st.dataframe()} \cite{streamlit_official}.

Streamlit tương thích với hầu hết các thư viện khoa học dữ liệu phổ biến bao gồm Pandas, NumPy, Matplotlib, Plotly, Scikit-learn, TensorFlow, PyTorch, Bokeh, Altair, OpenCV và Vega-Lite \cite{streamlit_official}. Ngoài ra, Streamlit Components cho phép cộng đồng xây dựng và chia sẻ các thành phần tùy chỉnh mở rộng khả năng của framework.

Ưu điểm: tốc độ phát triển nhanh, API đơn giản, không yêu cầu giao diện người dùng (front-end), tải lại tức thì (hot reload) tự động, tương thích rộng với các thư viện Python data science, triển khai dễ dàng \cite{streamlit_official}.

Nhược điểm: không phù hợp với các ứng dụng cần giao diện người dùng phức tạp hoặc tùy chỉnh cao, framework thiên về ứng dụng dữ liệu (data app) chứ không phải web framework đa năng (general-purpose), và các component mặc định có giới hạn về mặt thiết kế \cite{streamlit_official}.


\section{Docker}
Docker là nền tảng mã nguồn mở cho phép đóng gói ứng dụng và các phụ thuộc của nó vào container - một đơn vị chuẩn hóa chứa mọi thứ cần thiết để chạy phần mềm \cite{docker_official,docker_wikipedia}. Container được phát triển từ công nghệ cgroups và namespace của nhân Linux, cho phép cô lập tiến trình (process isolation) mà không cần máy ảo (virtual machine) riêng biệt \cite{docker_primer}. Với container, ứng dụng được cô lập hoàn toàn: chỉ cần cài đặt Docker Engine là có thể chạy bất kỳ container nào, bất kể hệ điều hành bên dưới \cite{docker_primer,docker_wikipedia}.

\subsection{Kiến trúc cơ bản của Docker}
Kiến trúc của Docker gồm ba thành phần chính \cite{docker_official}:

\begin{itemize}
    \item \textbf{Docker Daemon (dockerd):} Tiến trình nền (background process) chạy trên máy chủ, chịu trách nhiệm quản lý vòng đời container (tạo, chạy, dừng), xây dựng image, quản lý network và volume. Docker Daemon lắng nghe các API request từ Docker Client.

    \item \textbf{Docker Client:} Giao diện dòng lệnh (CLI) mà người dùng tương tác, ví dụ các lệnh docker build, docker run, docker pull. Docker Client giao tiếp với Daemon thông qua REST API \cite{docker_official}.

    \item \textbf{Registry (kho chứa image):} Nơi lưu trữ và phân phối Docker image. Docker Hub là registry công cộng mặc định, chứa hàng nghìn image sẵn có như ubuntu, python, postgres. Người dùng cũng có thể tự xây dựng private registry \cite{docker_official,docker_primer}.
\end{itemize}

\subsection{Image và Container}
Docker Image là một khuôn mẫu chỉ đọc (read-only template) chứa hệ điều hành gốc, mã ứng dụng, thư viện, và các cấu hình cần thiết. Image được xây dựng từ Dockerfile - tệp cấu hình khai báo từng bước để tạo nên image \cite{docker_official,docker_primer}.

Docker Container là một instance đang chạy (running instance) của một image. Nhiều container có thể chạy từ cùng một image mà không ảnh hưởng lẫn nhau, vì mỗi container có filesystem, network và process space riêng biệt \cite{docker_official}. Container được tạo bằng lệnh docker run; có thể dừng bằng docker stop và xóa bằng docker rm. Lệnh docker ps liệt kê các container đang chạy \cite{docker_wikipedia}.

Cơ chế layered filesystem cho phép các image chia sẻ chung các lớp (layer) dữ liệu, giảm đáng kể dung lượng lưu trữ và thời gian build. Khi một container thay đổi file, Docker sử dụng kỹ thuật sao chép khi ghi (copy-on-write): file được sao chép vào writable layer của container, tách biệt khỏi image gốc \cite{docker_primer}.

\subsection{Dockerfile và quy trình đóng gói}
Dockerfile là tệp văn bản chứa các instruction để tự động xây dựng image. Các instruction quan trọng bao gồm: FROM (image gốc), RUN (chạy lệnh trong quá trình build), COPY/ADD (sao chép tệp vào image), CMD/ENTRYPOINT (lệnh chạy khi container khởi động), EXPOSE (khai báo port), ENV (thiết lập biến môi trường) \cite{docker_official,docker_primer}.

Docker image được build theo từng lớp, mỗi instruction tạo ra một lớp mới. Docker cache các lớp này; khi thay đổi Dockerfile, chỉ các lớp phía sau thay đổi mới được rebuild, giảm đáng kể thời gian build \cite{docker_official}.

\subsection{Docker Compose}
Docker Compose là công cụ cho phép định nghĩa và chạy các ứng dụng đa container (multi-container) bằng một tệp cấu hình YAML duy nhất \cite{docker_official,docker_primer}. Thay vì chạy nhiều lệnh docker run riêng lẻ, Docker Compose cho phép khai báo tất cả container, network, và volume trong một tệp docker-compose.yml, sau đó khởi động hoặc dừng toàn bộ hệ thống bằng một lệnh duy nhất (docker compose up / docker compose down) \cite{docker_official}.

Ưu điểm của Docker: tính nhất quán giữa các môi trường (phát triển, kiểm thử, triển khai), cô lập hoàn toàn giữa các ứng dụng, khởi động nhanh (tính bằng giây so với phút của máy ảo), sử dụng tài nguyên hiệu quả hơn máy ảo nhờ dùng chung nhân, và hệ sinh thái phong phú với hàng nghìn image có sẵn trên Docker Hub \cite{docker_official,docker_primer}.

Nhược điểm của Docker: container dùng chung nhân của hệ điều hành chủ nên có bề mặt bảo mật chung, không phù hợp với các ứng dụng cần độc lập mạnh (như cơ sở dữ liệu multi-tenant), và yêu cầu quản lý cẩn thận về mạng và lưu trữ để tránh xung đột giữa các container \cite{docker_wikipedia}.


\section{Phân tích lựa chọn công nghệ}
Trong hệ thống gợi ý thời gian thực, ba nhóm nhu cầu chính được xác định: truyền tải và xử lý luồng sự kiện liên tục, huấn luyện và chạy mô hình học máy hoặc học sâu trên dữ liệu lớn, và phục vụ kết quả gợi ý với độ trễ thấp. Phần này phân tích và so sánh các công nghệ mã nguồn mở phù hợp cho đề tài. Nhóm lựa chọn những công cụ, công nghệ có mã nguồn mở hoặc có cộng đồng người dùng lớn, tài liệu chi tiết hoặc đã được chứng minh trong thực tế là hiệu quả và được sử dụng rộng rãi, tất nhiên là thêm những ưu thế về mặt kỹ thuật so với các công cụ tương tự, ưu thế trong việc áp dụng vào một dạng đề tài cụ thể


\subsection{Vai trò của từng công nghệ trong hệ gợi ý}

Công nghệ xử lý luồng sự kiện (event streaming) đóng vai trò tiếp nhận các tương tác người dùng (click, view, purchase) dưới dạng sự kiện, lưu trữ tạm thời và phân phối tới các pipeline xử lý. Khả năng replay events (đọc lại sự kiện cũ) là yếu tố then chốt cho phép phục hồi hoặc xử lý lại khi cần.

Công nghệ xử lý dữ liệu phân tán đảm nhận hai nhiệm vụ: batch processing (xử lý theo lô - huấn luyện mô hình trên toàn bộ dữ liệu lịch sử, tính toán , thống kê) và stream processing (cập nhật mô hình và kết quả gợi ý theo thời gian thực).

Lưu trữ cache phục vụ hai mục đích: làm buffer (bộ đệm) lưu lịch sử tương tác của phiên hiện tại (phục vụ tính toán theo phiên) và caching (làm bộ nhớ đệm) kết quả gợi ý để tránh tính toán lại.

Cơ sở dữ liệu quan hệ đóng vai trò kho dữ liệu tổng hợp, lưu trữ kết quả thống kê, log sự kiện, và các thông tin phục vụ dashboard phân tích.

\subsection{So sánh các công nghệ xử lý luồng sự kiện}

\begin{table}[H]
\centering
\resizebox{\textwidth}{!}{
\begin{tabular}{|l|c|c|c|}
\hline
\textbf{} & \textbf{Apache Kafka} & \textbf{RabbitMQ} & \textbf{Apache Pulsar} \\
\hline
Mô hình & Event streaming & Message broker & Event streaming \\
\hline
Thông lượng & Hàng triệu msg/s & 4K–10K msg/s & Hàng triệu msg/s \\
\hline
Lưu trữ & Log-based & Xóa sau ACK & Log-based (BookKeeper) \\
\hline
Replay events & Có (qua offset) & Không & Có (qua cursor) \\
\hline
Tích hợp Spark & Native connector & Cần custom & Native connector \\
\hline
Mở rộng & Theo chiều ngang & Nhân bản queue & Phân tách tính toán và lưu trữ \\
\hline
Độ trễ (p99) & 10–50ms & 5–15ms & 5–20ms \\
\hline
\end{tabular}
}
\caption{So sánh Apache Kafka với các công nghệ mã nguồn mở thay thế}
\label{tab:kafka_compare}
\end{table}

RabbitMQ phù hợp với việc gửi tin nhắn truyền thống ở thông lượng thấp, tuy nhiên không hỗ trợ việc lưu lại log nên không thể replay events (chạy lại sự kiện, gửi lại dữ liệu) - yếu tố quan trọng trong xử lý dữ liệu dạng luồng. Apache Pulsar có kiến trúc hiện đại với phân tách compute/storage nhưng cộng đồng nhỏ hơn, ecosystem chưa phong phú bằng Kafka \cite{kafka_vs_rabbitmq}. Kafka được lựa chọn vì lưu trữ log-based cho phép replay events phục vụ reprocessing và recovery; partition hỗ trợ xử lý song song; native Kafka source connector cho Spark Structured Streaming; và persistence cùng nhân bản (replication) đảm bảo khả năng chịu lỗi, cộng thêm với việc có cộng đồng lớn hỗ trợ và đã được chứng minh hiệu quả trong thực tế trong thời gian dài.

\subsection{So sánh các công nghệ xử lý dữ liệu phân tán}

\begin{table}[H]
\centering
\resizebox{\textwidth}{!}{
\begin{tabular}{|l|c|c|c|}
\hline
\textbf{} & \textbf{Apache Spark} & \textbf{Apache Flink} & \textbf{Hadoop MR} \\
\hline
Mô hình xử lý & Micro-batch & True streaming & Batch-only \\
\hline
Độ trễ streaming & 100ms–vài s & Sub-ms & Không áp dụng \\
\hline
Batch & In-memory & Bounded stream & Disk-based \\
\hline
Unified engine & Có & Có (stream-first) & Không \\
\hline
Tích hợp Kafka & Native & Native & Không \\
\hline
\end{tabular}
}
\caption{So sánh Apache Spark với Apache Flink và Hadoop MapReduce}
\label{tab:spark_compare}
\end{table}

Spark vượt trội về độ trễ streaming nhờ kiến trúc streaming thuần túy. Hadoop MapReduce dựa trên xử lý dựa trên đĩa gây độ trễ I/O (vào ra) đáng kể và không phù hợp cho trường hợp sử dụng thời gian thực \cite{spark_vs_flink_aws,spark_vs_flink_estuary}. Spark được lựa chọn vì là một công cụ xử lý được cả batch lẫn streaming trong một framework duy nhất; tích hợp Kafka gốc (native) qua Structured Streaming; xử lý trong bộ nhớ nhanh hơn MapReduce 10–100 lần; ngoài ra PySpark cho phép viết pipeline bằng Python.
\subsection{So sánh các công nghệ lưu trữ in-memory ( trong bộ nhớ hay cache)}

\begin{table}[H]
\centering
\begin{tabular}{|l|c|c|}
\hline
\textbf{} & \textbf{Redis} & \textbf{Memcached} \\
\hline
Kiểu dữ liệu & String, List, Hash, Set, ZSet, Stream & String only \\
\hline
Persistence & RDB(snapshot) + AOF(append only file) & Không \\
\hline
Replication & Master-replica (tích hợp sẵn) & Client-side sharding \\
\hline
TTL per-key & Có & Có \\
\hline
Độ trễ (p50) & 0.1–0.3ms & 0.08–0.2ms \\
\hline
Distributed & Redis Cluster & Client-side \\
\hline
Pub/Sub & Có & Không \\
\hline
\end{tabular}
\caption{So sánh Redis với Memcached}
\label{tab:redis_compare}
\end{table}

Memcached nhanh hơn 10–15\% cho những lệnh GET/SET cơ bản nhờ đa luồng, nhưng chỉ hỗ trợ kiểu dữ liệu string - không đủ cho các cấu trúc phức tạp như List cần cho lưu trữ lịch sử phiên \cite{redis_vs_memcached_techplained}.

Redis được lựa chọn vì kiểu dữ liệu List phù hợp cho việc lưu trữ session history; TTL per-key tự động xóa phiên hết hạn; độ trễ thấp phù hợp cho việc gợi ý thời gian thực; nhân bản và phân cụm cho khả năng mở rộng. Ngoài ra, Redis cũng đã được sử dụng nhiều và đều đã chứng minh được hiệu quả của mình, cũng như có cộng đồng hỗ trợ lớn và nhiều code mẫu, gợi ý.

\subsection{So sánh các công nghệ cơ sở dữ liệu quan hệ}

\begin{table}[H]
\centering
\resizebox{\textwidth}{!}{
\begin{tabular}{|l|c|c|c|}
\hline
\textbf{} & \textbf{PostgreSQL} & \textbf{MongoDB} & \textbf{MySQL} \\
\hline
Mô hình & Relational & Document (NoSQL) & Relational \\
\hline
ACID & Đầy đủ (SSI) & Multi-document (v4.0+) & Đầy đủ (InnoDB) \\
\hline
JSON support & JSONB + GIN & Native BSON & Hạn chế \\
\hline
Complex queries (truy vấn phức tạp) & Window functions, CTEs & Aggregation pipeline & Hạn chế hơn \\
\hline
Joins & Native, optimizer mạnh & \$lookup (hạn chế) & Native, optimizer yếu hơn \\
\hline
JDBC (Spark sink) & Có & Community driver & Có \\
\hline
Extensions & 1.000+ & Không & Hạn chế \\
\hline
Tuân thủ chuẩn SQL & Rất cao & Không (MQL) & Trung bình \\
\hline
\end{tabular}
}
\caption{So sánh PostgreSQL với MongoDB và MySQL}
\label{tab:postgresql_compare}
\end{table}

MongoDB phù hợp dữ liệu phi cấu trúc (do MongoDB là NoSQL, không phải là cơ sở dữ liệu quan hệ) nhưng yếu về những câu lệnh join phức tạp và việc tính toán, cộng gộp trong pipeline không mạnh bằng SQL window functions và CTEs. MySQL phổ biến nhưng query optimizer kém hơn cho những truy vấn phức tạp, JSON support hạn chế hơn JSONB cũng như ít tính năng cho các câu lệnh phân tích hơn PostgreSQL\cite{postgresql_vs_mongodb_devto,postgresql_vs_mongodb_modern,postgresql_vs_mysql_askantech}. PostgreSQL được lựa chọn vì MVCC hỗ trợ đồng thời đọc/ghi (concurrent read/write) hiệu quả; JSONB cho phép linh hoạt lưu trữ dữ liệu linh hoạt dạng Json nhưng lưu dưới dạng bit, không phải dưới dạng chuỗi do dạng chuỗi thì phải mã hóa, dạng bit thì tiết kiệm được bộ nhớ. Thêm một ưu thế của PostgreSQL là cộng đồng mã nguồn mở lớn và tốc độ phát triển tốt hơn nhiều so với MySQL, MySQL có bản mã nguồn mở được phát triển tiếp là MariaDB nhưng không được như PostgreSQL.

\end{document}
